% Sección 10.9: Representación de Funciones como Series de Taylor y Maclaurin

\section*{SECCIÓN 10.9 \quad Representación de Funciones como Series de Taylor y Maclaurin}

\noindent\textit{En los ejercicios siguientes, supóngase que todas las funciones tienen desarrollos en series de potencias.}

\vspace{0.3cm}
\noindent\textbf{En los ejercicios 1-12 use el método directo de los ejemplos 3, 4 y 7 para encontrar la serie de Taylor de $f$ en el valor dado de $c$. Encuentre también el radio de convergencia asociado.}
\begin{enumerate}
    \subimport{ejercicio_01/}{ejercicio.tex}
\subimport{ejercicio_02/}{ejercicio.tex}
\subimport{ejercicio_03/}{ejercicio.tex}
\subimport{ejercicio_04/}{ejercicio.tex}
\subimport{ejercicio_05/}{ejercicio.tex}
\subimport{ejercicio_06/}{ejercicio.tex}
\subimport{ejercicio_07/}{ejercicio.tex}
\subimport{ejercicio_08/}{ejercicio.tex}
\subimport{ejercicio_09/}{ejercicio.tex}
\subimport{ejercicio_10/}{ejercicio.tex}
\subimport{ejercicio_11/}{ejercicio.tex}
\subimport{ejercicio_12/}{ejercicio.tex}
\end{enumerate}

\vspace{0.3cm}
\noindent\textbf{En los ejercicios 13-24 encuentre la serie de Maclaurin de $f$ y su radio de convergencia usando una serie geométrica o bien derivando o integrando una serie geométrica.}
\begin{enumerate}
    \setcounter{enumi}{12}
    \subimport{ejercicio_13/}{ejercicio.tex}
\subimport{ejercicio_14/}{ejercicio.tex}
\subimport{ejercicio_15/}{ejercicio.tex}
\subimport{ejercicio_16/}{ejercicio.tex}
\subimport{ejercicio_17/}{ejercicio.tex}
\subimport{ejercicio_18/}{ejercicio.tex}
\subimport{ejercicio_19/}{ejercicio.tex}
\subimport{ejercicio_20/}{ejercicio.tex}
\subimport{ejercicio_21/}{ejercicio.tex}
\subimport{ejercicio_22/}{ejercicio.tex}
\subimport{ejercicio_23/}{ejercicio.tex}
\subimport{ejercicio_24/}{ejercicio.tex}
\end{enumerate}

\vspace{0.3cm}
\noindent\textbf{En los ejercicios 25-42 utilice cualquier método para encontrar la serie de Maclaurin de $f$ y su radio de convergencia. En particular, se puede hacer uso de las series conocidas 10.47, 10.49 y 10.50.}
\begin{enumerate}
    \setcounter{enumi}{24}
    \subimport{ejercicio_25/}{ejercicio.tex}
\subimport{ejercicio_26/}{ejercicio.tex}
\subimport{ejercicio_27/}{ejercicio.tex}
\subimport{ejercicio_28/}{ejercicio.tex}
\subimport{ejercicio_29/}{ejercicio.tex}
\subimport{ejercicio_30/}{ejercicio.tex}
\subimport{ejercicio_31/}{ejercicio.tex}
\subimport{ejercicio_32/}{ejercicio.tex}
\subimport{ejercicio_33/}{ejercicio.tex}
\subimport{ejercicio_34/}{ejercicio.tex}
\subimport{ejercicio_35/}{ejercicio.tex}
\subimport{ejercicio_36/}{ejercicio.tex}
\subimport{ejercicio_37/}{ejercicio.tex}
\subimport{ejercicio_38/}{ejercicio.tex}
\subimport{ejercicio_39/}{ejercicio.tex}
\subimport{ejercicio_40/}{ejercicio.tex}
\subimport{ejercicio_41/}{ejercicio.tex}
\subimport{ejercicio_42/}{ejercicio.tex}
\end{enumerate}

\vspace{0.3cm}
\begin{enumerate}
    \setcounter{enumi}{42}
    \subimport{ejercicio_43/}{ejercicio.tex}
\subimport{ejercicio_44/}{ejercicio.tex}
\end{enumerate}

\vspace{0.3cm}
\noindent\textbf{En los ejercicios 45-50 calcule la integral indefinida dada como serie infinita.}
\begin{enumerate}
    \setcounter{enumi}{44}
    \subimport{ejercicio_45/}{ejercicio.tex}
\subimport{ejercicio_46/}{ejercicio.tex}
\subimport{ejercicio_47/}{ejercicio.tex}
\subimport{ejercicio_48/}{ejercicio.tex}
\subimport{ejercicio_49/}{ejercicio.tex}
\subimport{ejercicio_50/}{ejercicio.tex}
\end{enumerate}

\vspace{0.3cm}
\noindent\textbf{En los ejercicios 51-56 utilice una serie para aproximar la integral definida dada con la precisión indicada.}
\begin{enumerate}
    \setcounter{enumi}{50}
    \subimport{ejercicio_51/}{ejercicio.tex}
\subimport{ejercicio_52/}{ejercicio.tex}
\subimport{ejercicio_53/}{ejercicio.tex}
\subimport{ejercicio_54/}{ejercicio.tex}
\subimport{ejercicio_55/}{ejercicio.tex}
\subimport{ejercicio_56/}{ejercicio.tex}
\end{enumerate}

\vspace{0.3cm}
\begin{enumerate}
    \setcounter{enumi}{56}
    \subimport{ejercicio_57/}{ejercicio.tex}
\subimport{ejercicio_58/}{ejercicio.tex}
\end{enumerate}

\vspace{0.3cm}
\noindent\textbf{En los ejercicios 59-64 encuentre la suma de la serie.}
\begin{enumerate}
    \setcounter{enumi}{58}
    \subimport{ejercicio_59/}{ejercicio.tex}
\subimport{ejercicio_60/}{ejercicio.tex}
\subimport{ejercicio_61/}{ejercicio.tex}
\subimport{ejercicio_62/}{ejercicio.tex}
\subimport{ejercicio_63/}{ejercicio.tex}
\subimport{ejercicio_64/}{ejercicio.tex}
\end{enumerate}

\vspace{0.3cm}
\noindent\textbf{En los ejercicios 65-68 aplique la multiplicación o división de series de potencias para encontrar los tres primeros términos no nulos de la serie de Maclaurin de la función dada.}
\begin{enumerate}
    \setcounter{enumi}{64}
    \subimport{ejercicio_65/}{ejercicio.tex}
\subimport{ejercicio_66/}{ejercicio.tex}
\subimport{ejercicio_67/}{ejercicio.tex}
\subimport{ejercicio_68/}{ejercicio.tex}
\end{enumerate}

\vspace{0.3cm}
\begin{enumerate}
    \setcounter{enumi}{68}
    \subimport{ejercicio_69/}{ejercicio.tex}
\subimport{ejercicio_70/}{ejercicio.tex}
\subimport{ejercicio_71/}{ejercicio.tex}
\subimport{ejercicio_72/}{ejercicio.tex}
\end{enumerate}
